% GENERATED FILE. Edit canonical lessons, then run npm run generate:books.
% canonical-source-hash: fnv1a64:1dec10a15568c67e
% canonical-lessons: SA-C390-laksanam, SA-C390-pratyasa, SA-C390-vyarthah, SA-C390-sarthakah, SA-C390-anubhavati

\chapter{Signs of What Is Coming}
\label{ch:sa-signs-of-what-is-coming}
\hlchaptermodality{390}

\begin{chapteropening}
\textbf{By the end of this chapter:} \emph{I can say what a sign tells us will happen.}

Complete the last lesson of chapter 390: I can say what a sign tells us will happen.
\end{chapteropening}

\section[\texorpdfstring{\sk{लक्षणम्} (lakṣaṇam)}{लक्षणम् (lakṣaṇam)}]{\sk{लक्षणम्} (lakṣaṇam) — a sign of something to come; a mark}

\label{lesson:SA-C390-laksanam}



\begin{quote}
Before the new one: say the Sanskrit for imagination, then the Sanskrit for true to fact.
\end{quote}

\subsection*{You'll want to know: \sk{लक्षणम्}}
\textbf{\sk{लक्षणम्}} — \emph{lakṣaṇam} — \textquotedblleft{}a sign of something to come; a mark\textquotedblright{}.

\begin{quote}
\sk{यदि} \sk{ज्वरस्य} \sk{लक्षणम्} \sk{अस्ति}, \sk{तर्हि} \sk{वैद्यं} \sk{पश्यतु।}
\end{quote}

\begin{quote}
\emph{yadi jvarasya lakṣaṇam asti, tarhi vaidyaṃ paśyatu.}
\end{quote}

\begin{quote}
If there is a sign of fever, see a doctor.
\end{quote}

A neuter noun.

\begin{cousinweb}
\textbf{\sk{लक्षणम्}} (\emph{lakṣaṇam}) comes from the root √lakṣ, \textquotedblleft{}to mark, to perceive\textquotedblright{}; MW gives \emph{lakṣaṇa} \textquotedblleft{}a mark, a sign, a symbol, a token, a characteristic\textquotedblright{} (MW).
\end{cousinweb}

\begin{grammarlens}[title={What you've built}]
One more word for saying what happens if.
\end{grammarlens}

\subsection*{Your turn}
\begin{itemize}
\raggedright
  \item \emph{Say it:} \emph{lakṣaṇam}
  \item \emph{Say it:} \emph{lakṣaṇam}, once more
  \item \emph{Say it:} \emph{yadi jvarasya lakṣaṇam asti, tarhi vaidyaṃ paśyatu.} — If there is a sign of fever, see a doctor.
  \item \emph{Recall:} say the Sanskrit for imagination, then the Sanskrit for true to fact, then say \emph{lakṣaṇam} again
\end{itemize}

\begin{tcolorbox}[breakable,colback=teal!4,colframe=teal!35!black,title={Before you move on}]
What does \sk{लक्षणम्} mean? (\textquotedblleft{}A sign of something to come; a mark\textquotedblright{}.) Say it once more.
\end{tcolorbox}
\section[\texorpdfstring{\sk{प्रत्याशा} (pratyāśā)}{प्रत्याशा (pratyāśā)}]{\sk{प्रत्याशा} (pratyāśā) — an expectation; a hope}

\label{lesson:SA-C390-pratyasa}



\begin{quote}
Before the new one: say the Sanskrit for true to fact, then the Sanskrit for a sign of something to come.
\end{quote}

\subsection*{You'll want to know: \sk{प्रत्याशा}}
\textbf{\sk{प्रत्याशा}} — \emph{pratyāśā} — \textquotedblleft{}an expectation; a hope\textquotedblright{}.

\begin{quote}
\sk{यदि} \sk{प्रत्याशा} \sk{अधिका}, \sk{तर्हि} \sk{दुःखम्} \sk{अपि} \sk{अधिकं} \sk{भवति।}
\end{quote}

\begin{quote}
\emph{yadi pratyāśā adhikā, tarhi duḥkham api adhikaṃ bhavati.}
\end{quote}

\begin{quote}
If expectation is great, sorrow is great too.
\end{quote}

A feminine noun.

\begin{cousinweb}
\textbf{\sk{प्रत्याशा}} (\emph{pratyāśā}) is built from \emph{prati}, \textquotedblleft{}towards\textquotedblright{}, and \emph{āśā}, \textquotedblleft{}hope\textquotedblright{}; MW gives \emph{praty-āśā} \textquotedblleft{}confidence, trust, hope, expectation\textquotedblright{} (MW).
\end{cousinweb}

\begin{grammarlens}[title={What you've built}]
One more word for saying what happens if.
\end{grammarlens}

\subsection*{Your turn}
\begin{itemize}
\raggedright
  \item \emph{Say it:} \emph{pratyāśā}
  \item \emph{Say it:} \emph{pratyāśā}, once more
  \item \emph{Say it:} \emph{yadi pratyāśā adhikā, tarhi duḥkham api adhikaṃ bhavati.} — If expectation is great, sorrow is great too.
  \item \emph{Recall:} say the Sanskrit for true to fact, then the Sanskrit for a sign of something to come, then say \emph{pratyāśā} again
\end{itemize}

\begin{tcolorbox}[breakable,colback=teal!4,colframe=teal!35!black,title={Before you move on}]
What does \sk{प्रत्याशा} mean? (\textquotedblleft{}An expectation; a hope\textquotedblright{}.) Say it once more.
\end{tcolorbox}
\section[\texorpdfstring{\sk{व्यर्थः} (vyarthaḥ)}{व्यर्थः (vyarthaḥ)}]{\sk{व्यर्थः} (vyarthaḥ) — useless; to no purpose}

\label{lesson:SA-C390-vyarthah}



\begin{quote}
Before the new one: say the Sanskrit for a sign of something to come, then the Sanskrit for an expectation.
\end{quote}

\subsection*{You'll want to know: \sk{व्यर्थः}}
\textbf{\sk{व्यर्थः}} — \emph{vyarthaḥ} — \textquotedblleft{}useless; to no purpose\textquotedblright{}.

\begin{quote}
\sk{यदि} \sk{यन्त्रं} \sk{न} \sk{चलति}, \sk{तर्हि} \sk{तत्} \sk{व्यर्थम्।}
\end{quote}

\begin{quote}
\emph{yadi yantraṃ na calati, tarhi tat vyartham.}
\end{quote}

\begin{quote}
If the machine does not run, it is useless.
\end{quote}

\begin{cousinweb}
\textbf{\sk{व्यर्थः}} (\emph{vyarthaḥ}) is built from \emph{vi}, \textquotedblleft{}without\textquotedblright{}, and \emph{artha}, \textquotedblleft{}purpose, use\textquotedblright{}; MW gives \emph{vy-artha} \textquotedblleft{}useless, unavailing, unprofitable, vain\textquotedblright{} (MW).
\end{cousinweb}

\begin{grammarlens}[title={What you've built}]
One more word for saying what happens if.
\end{grammarlens}

\subsection*{Your turn}
\begin{itemize}
\raggedright
  \item \emph{Say it:} \emph{vyarthaḥ}
  \item \emph{Say it:} \emph{vyarthaḥ}, once more
  \item \emph{Say it:} \emph{yadi yantraṃ na calati, tarhi tat vyartham.} — If the machine does not run, it is useless.
  \item \emph{Recall:} say the Sanskrit for a sign of something to come, then the Sanskrit for an expectation, then say \emph{vyarthaḥ} again
\end{itemize}

\begin{tcolorbox}[breakable,colback=teal!4,colframe=teal!35!black,title={Before you move on}]
What does \sk{व्यर्थः} mean? (\textquotedblleft{}Useless; to no purpose\textquotedblright{}.) Say it once more.
\end{tcolorbox}
\section[\texorpdfstring{\sk{सार्थकः} (sārthakaḥ)}{सार्थकः (sārthakaḥ)}]{\sk{सार्थकः} (sārthakaḥ) — worthwhile; meaningful}

\label{lesson:SA-C390-sarthakah}



\begin{quote}
Before the new one: say the Sanskrit for an expectation, then the Sanskrit for useless.
\end{quote}

\subsection*{You'll want to know: \sk{सार्थकः}}
\textbf{\sk{सार्थकः}} — \emph{sārthakaḥ} — \textquotedblleft{}worthwhile; meaningful\textquotedblright{}.

\begin{quote}
\sk{यदि} \sk{प्रयत्नः} \sk{फलं} \sk{ददाति}, \sk{तर्हि} \sk{सः} \sk{सार्थकः।}
\end{quote}

\begin{quote}
\emph{yadi prayatnaḥ phalaṃ dadāti, tarhi saḥ sārthakaḥ.}
\end{quote}

\begin{quote}
If an effort bears fruit, it is worthwhile.
\end{quote}

\begin{cousinweb}
\textbf{\sk{सार्थकः}} (\emph{sārthakaḥ}) is built from \emph{sa}, \textquotedblleft{}with\textquotedblright{}, and \emph{artha}, \textquotedblleft{}meaning, use\textquotedblright{}; MW gives \emph{sārthaka} \textquotedblleft{}having meaning, significant\textquotedblright{} and \textquotedblleft{}serviceable, useful\textquotedblright{} (MW).
\end{cousinweb}

\begin{grammarlens}[title={What you've built}]
One more word for saying what happens if.
\end{grammarlens}

\subsection*{Your turn}
\begin{itemize}
\raggedright
  \item \emph{Say it:} \emph{sārthakaḥ}
  \item \emph{Say it:} \emph{sārthakaḥ}, once more
  \item \emph{Say it:} \emph{yadi prayatnaḥ phalaṃ dadāti, tarhi saḥ sārthakaḥ.} — If an effort bears fruit, it is worthwhile.
  \item \emph{Recall:} say the Sanskrit for an expectation, then the Sanskrit for useless, then say \emph{sārthakaḥ} again
\end{itemize}

\begin{tcolorbox}[breakable,colback=teal!4,colframe=teal!35!black,title={Before you move on}]
What does \sk{सार्थकः} mean? (\textquotedblleft{}Worthwhile; meaningful\textquotedblright{}.) Say it once more.
\end{tcolorbox}
\section[\texorpdfstring{\sk{अनुभवति} (anubhavati)}{अनुभवति (anubhavati)}]{\sk{अनुभवति} (anubhavati) — experiences; feels}

\label{lesson:SA-C390-anubhavati}



\begin{quote}
Before the new one: say the Sanskrit for useless, then the Sanskrit for worthwhile.
\end{quote}

\subsection*{You'll want to know: \sk{अनुभवति}}
\textbf{\sk{अनुभवति}} — \emph{anubhavati} — \textquotedblleft{}experiences; feels\textquotedblright{}.

\begin{quote}
\sk{यदि} \sk{भवान्} \sk{ध्यायति}, \sk{तर्हि} \sk{शान्तिम्} \sk{अनुभवति।}
\end{quote}

\begin{quote}
\emph{yadi bhavān dhyāyati, tarhi śāntim anubhavati.}
\end{quote}

\begin{quote}
If you meditate, you feel peace.
\end{quote}

Past (he …): \sk{सः} \sk{शान्तिम्} \sk{अनुभूतवान्}, \emph{saḥ śāntim anubhūtavān}.

\begin{cousinweb}
\textbf{\sk{अनुभवति}} (\emph{anubhavati}) is built on \emph{anu}, \textquotedblleft{}after\textquotedblright{}, and the root √bhū, \textquotedblleft{}to be\textquotedblright{}; MW gives \emph{anu-√bhū} \textquotedblleft{}to attain, to equal\textquotedblright{} and \emph{anu-bhūta} \textquotedblleft{}perceived, understood\textquotedblright{}, and \textquotedblleft{}to experience\textquotedblright{} grows from those senses (MW).
\end{cousinweb}

\begin{grammarlens}[title={What you've built}]
One more word for saying what happens if.
\end{grammarlens}

\subsection*{Your turn}
\begin{itemize}
\raggedright
  \item \emph{Say it:} \emph{anubhavati}
  \item \emph{Say it:} \emph{anubhavati}, once more
  \item \emph{Say it:} \emph{yadi bhavān dhyāyati, tarhi śāntim anubhavati.} — If you meditate, you feel peace.
  \item \emph{Recall:} say the Sanskrit for useless, then the Sanskrit for worthwhile, then say \emph{anubhavati} again
\end{itemize}

\begin{tcolorbox}[breakable,colback=teal!4,colframe=teal!35!black,title={Before you move on}]
What does \sk{अनुभवति} mean? (\textquotedblleft{}Experiences; feels\textquotedblright{}.) Say it once more.
\end{tcolorbox}
